\documentclass{article}
\usepackage[margin=2cm]{geometry}
\usepackage{amsmath}
\usepackage{graphicx}
\usepackage{booktabs}
\usepackage[utf8]{inputenc}
\usepackage{ragged2e}
\setlength{\emergencystretch}{3em}
\begin{document}
\sloppy

\noindent\hfill\includegraphics[height=1.8em]{assets/figure_header_logo.png}\par\noindent\hrule\vspace{4pt}

\subsection*{0.1BPS Equations}


We now use the vanishing of the fermionic variations to obtain BPS equations for the warp.
factor and the three non-zero scalars.


\subsection*{0.1.1Dilatino equation and projector}





\begin{equation}
\begin{array} { c c c c c } { 1 } & {. } & { } & {. } & { } & { } & { } & { } & { } \\ { \frac { 1 } { 2 } \gamma ^ { 5 } \sigma ^ { \prime } \varepsilon _ { A } = N _ { 0 } \varepsilon _ { A } + N _ { 3 } \gamma ^ { 7 } ( \sigma ^ { 3 } ) ^ { B } { } _ { A } \varepsilon _ { B } } & { } & { } & { } & { } & { } & { ( 1 ) } \\ {. } & {. } & {.. } & { } & { } & { }\end{array}
\end{equation}


This equation can be interpreted as a projection condition on the spinors ea. Consistency
of this projection condition then requires that


\begin{equation}
\begin{array} { c c c c c } { \sigma ^ { \prime } = 2 \eta \sqrt { N _ { 0 } ^ { 2 } + N _ { 3 } ^ { 2 } } } & { } & { } & { } & { } & { } & { } & { ( 2 ) } \\ { \mathrm { ~ r a n e } \qquad, \quad. \quad. \cdot \quad. \quad. \quad } & { ~ }\end{array}
\end{equation}


where n = 1. Plugging this BPS equation back into then yields a second form of the.
projection condition,


\begin{equation}
\begin{array} { r } { \gamma ^ { 5 } \varepsilon _ { A } = G _ { 0 } \varepsilon _ { A } - G _ { 3 } \gamma ^ { 7 } ( \sigma ^ { 3 } ) ^ { B } \varepsilon _ { B } } \\ { \mathrm { ~. ~ } \mathbf { I } _ { \mathrm { s u t o m a l i c e n ~ }, \quad \ell \mathrm { ~ i b l ~ o ~ a d k, o n ~ } \mathrm { ~ D P C } } _ { \mathrm { ~ a c o s a l i c a l i c a n ~ }, \quad \mathbf { I } _ { \mathrm { s u t o m a l}}}}\end{array}
\end{equation}


which is more useful in the derivation of the other BPS equations. In the above, we have
defined


\begin{equation}
G _ { \mathbf { n } } = \eta \frac { N _ { 0 } } { \sqrt { N _ { 0 } ^ { 2 } + N _ { 3 } ^ { 2 } } } \qquad \qquad \qquad G _ { \mathbf { 4 } } = - \eta \frac { N _ { \mathbf { 3 } } } { \sqrt { N _ { \mathbf { 6 } } ^ { 2 } + N _ { 3 } ^ { 2 } } } \qquad ( 4 )
\end{equation}


\subsection*{0.1.2Gravitino equation}


The analysis of the gravitino equation oA = 0 proceeds in exactly the same way as for the
Lorentzian case studied in . The procedure gives rise to a first-order equation for the warp
factor f and an algebraic constraint. To avoid excessive overlap with that paper, we simply
cite the result,.


\begin{equation}
f ^ { \prime } = 2 ( G _ { 0 } S _ { 0 } + G _ { 3 } S _ { 3 } ) ~ ~ ~ ~ ~ ~ ~ ~ e ^ { - 2 f } = 4 ( G _ { 0 } S _ { 0 } + G _ { 3 } S _ { 3 } ) ^ { 2 } - 4 ( S _ { 0 } ^ { 2 } + S _ { 3 } ^ { 2 } ) ~
\end{equation}


\subsection*{0.1.3Gaugino equations}


Finally, we turn toward the gaugino equation dA = 0. Again the analysis of this equation.
proceeds in an exactly analogous manner to the Lorentzian case . The result is.


\begin{equation}
\begin{array} { c c } { \cos \phi ^ { 4 } ( \phi ^ { 0 } ) ^ { \prime } = - ( G _ { 0 } M _ { 0 } + G _ { 3 } M _ { 3 } ) ~ }\end{array}
\end{equation}


The right-hand sides of both equations are real, and thus give rise to real solutions when
appropriate initial conditions are imposed.

\end{document}
